\begin{frame}{Limitations and future directions}



\begin{columns}[T,onlytextwidth]
\begin{column}{.48\textwidth}

\begin{block}{\faLock~Current limitations}
\begin{itemize}
    \item depends on student readiness
    \item scalability considerations
    \item domain-specific transferability
\end{itemize}
\end{block}
\end{column}

\begin{column}{.48\textwidth}
\begin{block}{\faScrewdriverWrench~Design refinements}
\begin{itemize}
    \item earlier preparation for professional tooling
    \item improved stakeholder interaction
    \item structured team-building to accelerate trust in large groups
    \item scaling across courses and partners
\end{itemize}
\end{block}
\end{column}
\end{columns}

\end{frame}
\note{
This is a design-oriented case -- not a controlled study.  

It depends on student readiness,  
and there is some initial overhead.  

Early in the semester,  
students often experience the workload as heavy.  
But later they consistently describe this  
as one of the most valuable parts of their studies.  

So there is a clear shift  
from short-term effort  
to long-term value.  

From a teaching perspective,  
this creates a very \emph{dynamic environment} --  
but it also requires hands-on interaction.  
So scaling needs to be considered carefully.  
At the current size, it is manageable.  

In terms of transferability,  
this fits naturally within BI --  
but the underlying design can extend to other courses.  
But it requires both students and instructors  
to commit to the structure.  

That is why we are now introducing  
tooling and practices earlier in the curriculum.  

And one ongoing refinement  
is strengthening the initial phase --  
giving students more time  
to build trust,  
understand the data,  
and align as a team  --
before introducing more advanced methods.
}